\documentclass[11pt]{article}
\usepackage[margin=0.9in]{geometry}
\usepackage{amsmath,amssymb,amsthm}
\usepackage[T1]{fontenc}
\usepackage{lmodern}
\usepackage{microtype}
\usepackage[colorlinks=true,linkcolor=blue,citecolor=blue,urlcolor=blue]{hyperref}
\newtheorem{theorem}{Theorem}
\newtheorem{lemma}[theorem]{Lemma}
\newcommand{\A}{\mathcal A}
\newcommand{\T}{\mathcal T}
\newcommand{\D}{\mathcal D}
\newcommand{\B}{\mathcal B}
\newcommand{\C}{\mathcal C}
\newcommand{\WOT}{\mathrm{WOT}}
\newcommand{\norm}[1]{\lVert#1\rVert}
\title{Tail generation for symmetric free-product states:\\a general deduction}
\author{Alec Kriebel\\
\small Independent researcher\\
\small \href{https://orcid.org/0009-0001-9320-500X}{ORCID 0009-0001-9320-500X}}
\date{22 October 2026 \quad Version 1.0.0}
\begin{document}
\maketitle
\begin{abstract}
For every unital \(C^*\)-algebra and every state on its countable universal
unital free product invariant under finite permutations, the recursive
tail \(C^*\)-algebra of Dykema--K\"ostler--Williams generates the tail von
Neumann algebra in the original GNS representation. More precisely,
the canonical expectation of the represented free-product algebra is
weak-operator dense in the tail with the exact norm bound of each target.
This is a short general deduction from the earlier coefficient-stabilization
method and Kaplansky density, answering the explicit generation question.
No faithfulness or separability is required. We explain the close relation
to Liu's bilateral future-expectation argument and retain bounded
priority and source-version qualifications.
\end{abstract}

\section{The question, its definition, and the conclusion}
Let \(A\) be a unital \(C^*\)-algebra, let
\(\A=*_{j\in\mathbb N}A\) be its universal unital free product, with
\(\mathbb N=\{1,2,\ldots\}\), and let \(\psi\) be a state invariant under
every finite permutation of the copies. Write \((\pi,H,\xi)\) for its
GNS representation, \(\pi_j:A\to B(H)\) for the represented \(j\)-th
copy, and
\[
 M=\pi(\A)'',\qquad
 \C_N=C^*(\pi_j(A):j>N),\qquad
 \T=\bigcap_{N\geq1}\C_N'',\qquad
 Q=C^*(\T\cup\pi(\A)).
\]
All closures and commutants here are in \(B(H)\), with its original
cyclic representation; \(C^*\)-algebras generated by unital copies
contain the common unit.

We use the canonical expectation \(F:Q\to\T\) of
Dykema--K\"ostler--Williams (DKW), Theorem~5.1.10(iv) \cite{DKW}.
The source provides an isometric shift \(\alpha\) on \(Q\), satisfying
\(\alpha(\pi_j(a))=\pi_{j+1}(a)\), and
\begin{equation}\label{shiftlimit}
 F(x)=\WOT\!-\!\lim_{k\to\infty}\alpha^k(x)\quad(x\in Q).
\end{equation}
It also provides spatial finite-permutation actions \(\beta_\sigma\)
with implementing unitaries fixing \(\xi\).
These are results for general symmetric states in DKW~\S5, rather than
assumptions of quantum symmetry. In particular \(F\) is contractive
and fixes \(\T\); a normal extension to \(M\) is not asserted.

Precisely as in DKW Definition~5.3.1, \(\D\) is the least unital
\(C^*\)-subalgebra of \(\T\) such that, for every \(m\geq1\), every
\(a_1,\ldots,a_m\in A\), every \(i_1,\ldots,i_m\in\mathbb N\), and
every \(d_1,\ldots,d_{m-1}\in\D\),
\[
 F\bigl(\pi_{i_1}(a_1)d_1\pi_{i_2}(a_2)\cdots
          d_{m-1}\pi_{i_m}(a_m)\bigr)\in\D .
\]
There is no distinct-index restriction. The family is nonempty because
\(\T\) has this property, and its intersection defines \(\D\).

\begin{theorem}\label{main}
For every \(A,\psi\) as above, \(\D''=\T\). In fact, for each \(t\in\T\)
there is a net \(x_\lambda\in\pi(\A)\) such that
\[
 \norm{x_\lambda}\leq\norm{t},\qquad
 \norm{F(x_\lambda)}\leq\norm{t},\qquad
 F(x_\lambda)\longrightarrow t\quad\hbox{in }\WOT .
\]
Consequently \(\overline{F(\pi(\A))}^{\,\WOT}=\T\).
\end{theorem}

This answers DKW Question~5.3.5 and the separately posed Question~2
in the 2015 and 2016 Oberwolfach reports \cite{OWR15,OWR16}.
Their other question, whether \(\D\subseteq\pi(\A)\), is separate and
is not answered here. DKW Theorem~7.3 already proves tail generation
for quantum symmetric states, including nonfaithful states.
The present statement imposes no faithfulness of \(\psi\) or \(\pi\),
separability, trace, factoriality, or quantum symmetry.

\section{Future compression and bounded approximation}
Put
\[
 \B_N=C^*(\pi_j(A):1\leq j\leq N),\qquad
 H_N=\overline{\B_N\xi},\qquad
 P_N:H\to H_N
\]
with \(P_N\) the orthogonal projection.
\begin{lemma}\label{compression}
For every \(N\geq1\) and \(x\in\C_N\),
\[
 P_NF(x)P_N=P_NxP_N .
\]
\end{lemma}
\begin{proof}
First take a finite-copy polynomial \(x\) whose indices belong to a finite
set \(S\subset\{N+1,N+2,\ldots\}\). Fix \(k\geq1\).
The injection \(i\mapsto i+k\) from \(S\) onto \(S+k\) extends to a
permutation of the finite set \(U=S\cup(S+k)\): complete it by any
bijection \(U\setminus S\to U\setminus(S+k)\). Extend that permutation
by the identity off \(U\). It fixes every first-\(N\) copy and transports
every occurrence of each index simultaneously. Overlap between \(S\)
and \(S+k\), repeated indices, and shared indices among different
monomials cause no problem.

For \(y,z\in\B_N\), state invariance gives, for every \(k\),
\[
 \langle\alpha^k(x)y\xi,z\xi\rangle
 =\psi\bigl(\beta_\sigma(z^*xy)\bigr)
 =\psi(z^*xy)=\langle xy\xi,z\xi\rangle.
\]
The notation \(\psi\) on represented elements means its GNS vector state.
The permutation action is well-defined even in a nonfaithful GNS
representation: state preservation preserves the GNS null space and
the GNS inner product, and its unitary fixes \(\xi\).
Taking \(k\to\infty\) uses \eqref{shiftlimit} only for this fixed input
\(x\). Future polynomials are norm dense in \(\C_N\); contractivity
of \(F\) extends the identity to every \(x\in\C_N\).
Finally \(\B_N\xi\) is dense in \(H_N\), and boundedness of \(F(x)\)
and \(x\) extends it to all vector pairs in \(H_N\).
\end{proof}

We recall exactly the norm-ball form of Kaplansky density used below
\cite{Kaplansky}. If \(C\subset B(H)\) is a unital concrete \(C^*\)-algebra
and \(t\in C''\), then its ball of radius \(\norm t\) is strongly, hence
weak-operator, dense at \(t\). For completeness, the reduction of the
nonselfadjoint case to the usual selfadjoint theorem preserves this
exact radius. Apply selfadjoint norm-ball Kaplansky in \(M_2(C)\) to
\[
 Y=\begin{pmatrix}0&t\\t^*&0\end{pmatrix}
       \in M_2(C''),\qquad \norm Y=\norm t .
\]
The upper-right corners of selfadjoint approximants of norm at most
\(\norm t\) lie in \(C\), have that same norm bound, and converge strongly
to \(t\), by testing vectors \((0,\eta)\). Thus no factor two or
norm-attaining lift from an abstract quotient is needed. The matrix device
already appears in Kaplansky's proof, which credits P.~R.~Halmos.

\begin{proof}[Proof of Theorem~\ref{main}]
Fix \(t\in\T\), and set \(r=\norm t\). If \(r=0\), take the constant
zero net. The increasing union \(\bigcup_N H_N\) is dense in \(H\):
every finite sum of copywords uses only finitely many indices, and
\(\xi\) is cyclic for \(\pi(\A)\). This does not require a countable
dense subset of \(A\) or \(H\). Hence \(P_N\to I\) strongly.

Specify a WOT neighborhood of \(t\) by finitely many vector pairs
\((\eta_l,\zeta_l)\), \(1\leq l\leq m\), and an error \(\varepsilon>0\).
Choose a single \(N\) so large that, for all \(l\),
\begin{equation}\label{vectors}
 2r\bigl(\norm{\eta_l-P_N\eta_l}\norm{\zeta_l}
       +\norm{P_N\eta_l}\norm{\zeta_l-P_N\zeta_l}\bigr)
       <\varepsilon/2 .
\end{equation}
Since \(t\in\C_N''\), the preceding norm-ball theorem gives
\(x\in\C_N\), \(\norm x\leq r\), with
\[
 |\langle(x-t)P_N\eta_l,P_N\zeta_l\rangle|<\varepsilon/2
       \quad(1\leq l\leq m).
\]
The algebra \(\C_N\) is contained in \(\pi(\A)\), since the image of a
\(C^*\)-homomorphism is norm closed and contains every copy image.

Put \(d=F(x)\). Contractivity gives \(\norm d\leq r\) and
\(\norm{d-t}\leq2r\). Lemma~\ref{compression} makes the compressed
coefficient of \(d-t\) equal to the displayed coefficient of \(x-t\).
Splitting the original vectors at \(P_N\) therefore gives
\begin{align*}
 |\langle(d-t)\eta_l,\zeta_l\rangle|
 &\leq|\langle(d-t)P_N\eta_l,P_N\zeta_l\rangle|\\
 &\quad+2r\bigl(\norm{\eta_l-P_N\eta_l}\norm{\zeta_l}
       +\norm{P_N\eta_l}\norm{\zeta_l-P_N\zeta_l}\bigr)
 <\varepsilon .
\end{align*}
Directing WOT neighborhoods by reverse inclusion gives the required
bounded net of represented inputs. Conversely its outputs lie in the
WOT-closed algebra \(\T\).

By the definition of \(\D\), setting every inserted \(d_j=1\) shows
that \(F\) of each finite copyword belongs to \(\D\).
Their linear span is norm dense in \(\pi(\A)\). Boundedness of \(F\)
and norm closure of \(\D\) imply \(F(\pi(\A))\subseteq\D\subseteq\T\).
The established density now gives
\(\T\subseteq\D''\subseteq\T\), as required.
\end{proof}

\section{What the argument uses, and its relation to prior work}
The proof never passes a WOT-convergent input net through \(F\).
For each chosen future algebra it uses an exact compression identity,
then controls the remaining vector coefficients by a uniform norm bound.
The uniform bound is essential; convergence on a dense vector subspace
alone need not imply WOT convergence.

Nonnormality is a real obstruction to a different shortcut.
In the compact-plus-scalar example recorded by DKW, Example~5.2.1
(credited there to Liu), take \(H=\ell^2(\{0,1,\ldots\})\),
\(a_j=e_{0j}+e_{j0}\), \(Q=K(H)+\mathbb CI\), and
\(\T=\mathbb Ce_{00}+\mathbb C(I-e_{00})\). The canonical expectation is
\[
 F(\lambda I+K)=\lambda I+K_{00}e_{00}.
\]
The projections \(q_m=\sum_{j=0}^m e_{jj}\) increase strongly to \(I\),
but \(F(q_m)=e_{00}\ne F(I)=I\), excluding a normal extension fixing
these values. Nevertheless \(F(a_j^2)=e_{00}\), so ordinary moments
already span the tail in this example, as the theorem predicts.

\paragraph{Attribution and the narrow contribution.}
DKW Lemma~5.1.8, the proof of Theorem~5.1.10, and Remark~5.1.11 already
contain finite-word coefficient stabilization and bounded-vector
extension. Liu's bilateral treatment \cite[\S6.2]{Liu}, especially
Lemma~6.8 and Proposition~6.10 in the inspected arXiv version, gives
a closely related past/future identity and normality of the expectation
on each fixed future von Neumann algebra. Neither a new coefficient
principle nor a new Kaplansky method is claimed here. The contribution
is the short general deduction for the explicit DKW recursive
\(C^*\)-algebra in its original representation.

The relation to Liu is closer than a superficial distinction between
bilateral and unilateral indexing suggests. Extend the symmetric law
to indices \(\mathbb Z\); in its bilateral GNS representation the shift
is spatial and normal. The same prior stabilization argument, extended
to blocks of elements of arbitrary \(A\) and bounded nets, yields a
normal expectation from each fixed future algebra onto its positive
tail. Restriction to the positive cyclic subspace gives a normal
surjection \(\rho:M_+\to M\). Its image of the bilateral positive tail
is exactly \(\T\): for \(t\in\T\), the norm-\(\norm t\) lifts in the
successive future algebras form nonempty nested ultraweakly compact
sets, and their intersection provides a tail lift. Compatibility of
the shift limits on \(C^*\)-inputs then gives the same density conclusion.
One must prove this tail-lifting step; images cannot simply be commuted
with intersections. Nor need the bilateral expectation descend through
\(\rho\) on all of \(M_+\). The accompanying
\texttt{RELATION\_TO\_PRIOR\_WORK.md} supplies this bridge in detail.
It emphasizes the corollary character of the result, rather than
asserting that Liu stated the full arbitrary-\(A\) DKW answer.

\paragraph{Versions, priority limits, and provenance.}
The relevant DKW \S5 and Liu \S6.2 bodies were read in the current
responses at the versioned arXiv URLs below. The DKW response has
internal running headers dated August 2021 although arXiv labels v3
as September 2014; no identity with historical 2014 or journal bytes is
inferred. Both AMS full-body routes were inaccessible, and Liu's
dissertation body was not recovered. A bounded primary-source
investigation found no authenticated earlier publication of the full
general answer. It does not certify exclusive firstness, worldwide
novelty, or continuing openness. The proof was already public in project
\href{https://github.com/AlecKriebel/Math/pull/1329}{PR1329}
(13 October 2026); this note is dated 22 October 2026 and is not backdated.
Catalog numbers 30003208 and30002897 name the same generation question.

\paragraph{Verification, authorship, and status.}
The source archive includes portable standard-library exact controls
and deliberately false variants for the finite permutation and
compact-plus-scalar mechanisms. These finite checks supplement the
written proof and do not certify arbitrary \(C^*\)-algebras or priority.
AI tools were extensively involved in research planning, mathematical
exploration, proof development and reconstruction, computation, source
and literature investigation, writing, editing, and adversarial
verification. Alec Kriebel is the named human author.
This work is unrefereed, without conventional independent human peer
review, journal acceptance, or formal proof-assistant certification;
internal AI audits do not supply those certifications.
Prose is licensed CC BY 4.0; accompanying authored code is MIT licensed.

\begingroup
\small
\begin{thebibliography}{9}
\bibitem{DKW}
K.~J.~Dykema, C.~K\"ostler and J.~D.~Williams,
\emph{Quantum symmetric states on free product \(C^*\)-algebras},
Trans.\ Amer.\ Math.\ Soc.\ \textbf{369} (2017), 645--679.
\href{https://doi.org/10.1090/tran6661}{doi:10.1090/tran6661}.
Inspected: \href{https://arxiv.org/pdf/1305.7293v3}{arXiv:1305.7293v3},
\S\S5.1--5.3 and the quantum-symmetric theorem.
\bibitem{Liu}
W.~Liu, \emph{Extended de Finetti theorems for boolean independence and
monotone independence}, Trans.\ Amer.\ Math.\ Soc.\ \textbf{370}
(2018), 1959--2003.
\href{https://doi.org/10.1090/tran/7034}{doi:10.1090/tran/7034}.
Inspected numbering:
\href{https://arxiv.org/pdf/1505.02215v1}{arXiv:1505.02215v1}, \S6.2.
\bibitem{Kaplansky}
I.~Kaplansky, \emph{A theorem on rings of operators},
Pacific J.\ Math.\ \textbf{1} (1951), 227--232.
\href{https://msp.org/pjm/1951/1-2/pjm-v1-n2-p06-s.pdf}{Primary paper},
Theorem 1.
\bibitem{OWR15}
K.~Dykema (joint with C.~K\"ostler and J.~D.~Williams),
\emph{Tail algebras of symmetric states and de Finetti theorems},
Oberwolfach Report 28/2015, printed pp.1618--1619, Question 2.
\href{https://doi.org/10.4171/OWR/2015/28}{doi:10.4171/OWR/2015/28}.
\href{https://ems.press/content/serial-article-files/46575}{Primary report}.
\bibitem{OWR16}
K.~Dykema, \emph{Tail algebras, amalgamated free products and KMS states},
Oberwolfach Report 40/2016, printed pp.2334--2335, Question 2.
\href{https://doi.org/10.4171/OWR/2016/40}{doi:10.4171/OWR/2016/40}.
\href{https://ems.press/content/serial-article-files/46645}{Primary report}.
\end{thebibliography}
\endgroup
\end{document}
